\documentclass[11pt]{article}
\usepackage{amsmath,amssymb,amsthm}

\newtheorem{theorem}{Theorem}
\newtheorem{corollary}[theorem]{Corollary}
\newtheorem{lemma}[theorem]{Lemma}

\newcommand{\Up}{U_+}

\title{Decoder ladders and predecessor blockers at the record atom $1113$}
\author{}
\date{27 August 2026}

\begin{document}
\maketitle

Let
\[
 x_{t+1}=\Up(x_t):=\frac{3x_t+1}{2^{a_t}},
 \qquad a_t=v_2(3x_t+1),
 \tag{1}
\]
be an accelerated orbit of positive odd integers.  On an interior
deficit-four record tail, write
\[
 A:=1,\qquad B:=12,\qquad C:=1113,\qquad D:=1122
 \tag{2}
\]
for the four record atoms, and put $b(x):=v_3(x+1)$ at a record state.
Their decoder transport is
\[
 A:b\mapsto b+1,
 \qquad B:b\mapsto0,
 \qquad C:b\mapsto1,
 \qquad D:b\mapsto0.
 \tag{3}
\]
For a normalized unbounded ray, let $Q$ denote the coefficient at a
specified state and let $\mathcal K$ be its stationary height.  The
synchronized companion satisfies
\[
 2^{a_t}e_{t+1}=3e_t-1,
 \qquad e_t>1,
 \qquad x_t+e_t=\mathcal KQ_t.
 \tag{4}
\]

\begin{theorem}[the complete decoder ladder below a $C$ atom]
Suppose that a record transition starts at $x$, has atom $C=1113$, and
ends at $y$.  Put
\[
 b:=v_3(x+1).
 \tag{5}
\]
For every integer $s$ with $0\leq s\leq b$, define
\[
 \boxed{
 z_s+1:=\frac{2^{s-1}(x+1)}{3^s}.}
 \tag{6}
\]
At $s=0$, formula (6) means $z_0=(x-1)/2$.  Every $z_s$ is a positive
odd integer, its exact accelerated word to $y$ is
\[
 \boxed{1^s1121,}
 \tag{7}
\]
and hence
\[
 \Up^{s+4}(z_s)=y.
 \tag{8}
\]

If $e_x$ is the synchronized companion value at $x$, put
\[
 \epsilon_s
 :=1+\left(\frac23\right)^s\frac{e_x-1}{2}.
 \tag{9}
\]
Then $\epsilon_s>1$, it is transported to the original companion at $y$
by the word in (7), and
\[
 \boxed{
 z_s+\epsilon_s
 =\frac{2^{s-1}}{3^s}(x+e_x).}
 \tag{10}
\]

If $x$ lies on the orbit of a globally least positive odd unbounded seed
$n$, then, for every $0\leq s\leq b$,
\[
 \boxed{
 x+1\geq2(n+1)\left(\frac32\right)^s.}
 \tag{11}
\]
If $Q$ is the coefficient at $x$, then
\[
 \boxed{
 Q>\frac{2(n+1)}{\mathcal K}
       \left(\frac32\right)^s.}
 \tag{12}
\]
In particular, (11)--(12) hold with $s=b$.
\end{theorem}

\begin{proof}
The exact $C$ congruence and affine formula are
\[
 x\equiv79\pmod {128},
 \qquad
 y=\frac{81x+65}{64}.
 \tag{13}
\]
Thus $v_2(x+1)=4$.  The state $z_0=(x-1)/2$ has exact word $1121$ and
reaches $y$.

For $0\leq s\leq b$, equation (6) gives
\[
 z_s+1=\frac{2^s(z_0+1)}{3^s},
 \qquad
 v_2(z_s+1)=s+3.
 \tag{14}
\]
Under an exponent-one step, the shifted state $X+1$ is multiplied by
$3/2$.  Consequently the first $s$ steps from $z_s$ are exact
exponent-one steps and end at $z_0$.  Appending the exact word $1121$
proves (7)--(8).

Backward transport of a companion through one exponent-one step is the
map
\[
 E\longmapsto\frac{2E+1}{3},
 \qquad
 \frac{2E+1}{3}-1=\frac23(E-1).
 \tag{15}
\]
At $z_0$, the companion value is $(e_x+1)/2$.  Iterating (15) proves
(9), and addition of the state and companion formulas gives (10).

Every $z_s$ reaches the unbounded tail, so global leastness gives
$z_s\geq n$.  Equation (6) now proves (11).  Since
$x+e_x=\mathcal KQ$, equations (10), $z_s\geq n$, and
$\epsilon_s>1$ give
\[
 \frac{2^{s-1}}{3^s}\mathcal KQ
 =z_s+\epsilon_s>n+1,
\]
which is (12).
\end{proof}

\begin{corollary}[pullback through a one-atom run]
Suppose that a record state $p$ is followed by the atom block $A^rC$,
where $r\geq0$, and let $y$ be its terminal state.  Then
\[
 \boxed{h:=\frac{p-1}{2}}
 \tag{16}
\]
is a positive odd integer with exact word $1^r1121$ to $y$.  If $e_p$ is
the synchronized companion at $p$, the alternative companion at $h$ is
$(e_p+1)/2$, and
\[
 h+\frac{e_p+1}{2}=\frac{p+e_p}{2}.
 \tag{17}
\]
On a globally least unbounded ray, $p\geq2n+1$.
\end{corollary}

\begin{proof}
Let $x$ be the start of the terminal atom $C$.  The $r$ atoms $A$ give
\[
 x+1=\frac{3^r(p+1)}{2^r},
 \qquad
 b(x)=b(p)+r\geq r.
 \tag{18}
\]
Taking $s=r$ in (6) gives $z_r=(p-1)/2$.  The assertions now follow
from Theorem~1: indeed, forward companion transport through $A^r$ gives
$e_x-1=(3/2)^r(e_p-1)$, so (9) becomes
$\epsilon_r=(e_p+1)/2$.
\end{proof}

\begin{lemma}[least four-step predecessor of half a $C$ endpoint]
Suppose that a $C$ atom starts at $q$ and ends at $p$, and put
\[
 h:=\frac{p-1}{2}.
 \tag{19}
\]
If $h$ is odd, then every positive odd $w$ satisfying
$\Up^4(w)=h$ obeys
\[
 \boxed{w\geq4q-1.}
 \tag{20}
\]
Equality holds, and the unique exponent word attaining it is $1215$.
\end{lemma}

\begin{proof}
The affine equation for $C$ gives
\[
 p=\frac{81q+65}{64},
 \qquad
 h=\frac{81q+1}{128},
 \qquad
 h\equiv50\pmod {81}.
 \tag{21}
\]
In particular, $q\equiv79\pmod {128}$ and $q\geq79$.

Let $v=(c_1,c_2,c_3,c_4)$ be the exact word from $w$ to $h$, let
$B=c_1+c_2+c_3+c_4$, and put
\[
 C(v)=27+9\,2^{c_1}+3\,2^{c_1+c_2}
                  +2^{c_1+c_2+c_3}.
 \tag{22}
\]
Then
\[
 81w+C(v)=2^Bh.
 \tag{23}
\]
Writing
\[
 s_1=c_1,
 \qquad s_2=c_1+c_2,
 \qquad s_3=c_1+c_2+c_3,
\]
the integrality condition modulo $81$ is
\[
 27+9\,2^{s_1}+3\,2^{s_2}+2^{s_3}
 \equiv50\,2^B\pmod {81},
 \qquad
 1\leq s_1<s_2<s_3\leq B-1.
 \tag{24}
\]
Using the residues of $2^j$ modulo $81$, the finite check for
$4\leq B\leq9$ is
\[
\begin{array}{c|rrrrrr}
B&4&5&6&7&8&9\\ \hline
50\,2^B\bmod81&71&61&41&1&2&4\\
\#\text{ solutions of (24)}&0&0&0&0&0&1.
\end{array}
\tag{25}
\]
For $B=9$, the unique solution is
$(s_1,s_2,s_3)=(1,3,4)$, hence
\[
 v=1215,
 \qquad C(v)=85.
 \tag{26}
\]
Equations (21), (23), and (26) give
\[
 w=\frac{512h-85}{81}=4q-1.
 \tag{27}
\]
The five states for the proposed word are
\[
 4q-1,
 \quad 6q-1,
 \quad \frac{9q-1}{2},
 \quad \frac{27q-1}{4},
 \quad \frac{81q+1}{128}=h.
\]
The congruence $q\equiv79\pmod {128}$, together with the hypothesis that
$h$ is odd, shows that the extracted powers of two are exactly
$1,2,1,5$.

It remains to exclude a smaller predecessor with $B\geq10$.  Since every
$c_i\geq1$, equation (22) gives
\[
 C(v)
 \leq27+9\,2^{B-3}+3\,2^{B-2}+2^{B-1}
 =27+\frac{19}{8}2^B.
 \tag{28}
\]
As $h\geq50$ and $B\geq10$,
\[
\begin{aligned}
 &(2^B-512)h-(C(v)-85)\\
 &\quad\geq
 50(2^B-512)-\left(\frac{19}{8}2^B-58\right)>0.
\end{aligned}
\tag{29}
\]
Subtracting (27) from (23) now gives $w>4q-1$.  This proves (20) and the
equality statement.
\end{proof}

\begin{theorem}[nontrivial record atoms block the half-seed pullback]
Suppose that a nontrivial atom $P\in\{B,C,D\}$ ends at a record state
$p$, and that $p$ is followed by $A^rC$ for some $r\geq0$.  Let
$h=(p-1)/2$ be the exact coalescing seed from Corollary~2.
\begin{enumerate}
 \item If $P=B$ or $P=D$, then $3\mid h$.  Hence $h$ has no accelerated
 positive odd predecessor.
 \item If $P=C$ starts at $q$, then every four-step accelerated pullback
 of $h$ is at least $4q-1>q$.
\end{enumerate}
Thus the half-seed construction propagates backward through the complete
intervening run of atoms $A$, but across the preceding nontrivial atom it
either terminates or expands beyond that atom's starting state.
\end{theorem}

\begin{proof}
Write $Y=p+1$.  If $P=B$, its shifted-state equation has the form
\[
 8Y=9Y_0+4,
\]
and if $P=D$, it has the form
\[
 64Y=81Y_0+56.
\]
In both cases $Y\equiv2\pmod3$.  Therefore
\[
 h=\frac{Y}{2}-1\equiv0\pmod3.
\]
A multiple of $3$ has no accelerated odd predecessor, proving the first
part.  The second part is Lemma~3.
\end{proof}

\begin{corollary}[localized remaining gate in the infinite-$C$ branch]
On an interior record tail with infinitely many atoms $C$, every local
$C$ half-coalescence has the complete ladder (6)--(12).  Pulling it toward
the fixed initial seed by the preceding atoms $A$ is exact.  Apart from at
most the first $C$ after the chosen tail origin, the first preceding atom
in $\{B,C,D\}$ is a blocker in the sense of Theorem~4.
Consequently, any global descent obtained from infinitely many $C$ atoms
must use an inverse branch other than the direct half-seed ladder across
at least one nontrivial record atom.
\end{corollary}

\begin{proof}
Apply Corollary~2 to the maximal run $A^r$ preceding each selected atom
$C$, and then apply Theorem~4 to the preceding nontrivial atom.
\end{proof}

\end{document}
